\section{Related Work}

% early history
Ever since the first chatbots\footnote{The first recorded conversation between the chatbots ELIZA and PARRY is available here: \url{https://www.rfc-editor.org/rfc/rfc439}}, humans have been fascinated by the ability of computers to communicate in a human-like manner.
Recent advances in LLMs to reason and solve complex tasks \citep{OpenAIAAA24a} lead to a surge in studies on multi-agent systems \citep{ZhaoHXL23a, XuYLW23a, SuzgunK24a}.
\citet{GuoCWC24a} conduct a literature review on multi-agent LLMs.
Their taxonomy points to two main areas in which multi-agent LLMs are used: simulation and problem-solving.
Our investigation explores problem-solving because these tasks offer a controllable environment to probe components in multi-agent interaction \citep{YinSCG23a, DuLTT23}. %, we poswhich %is also the application we focus on in this work.
%They highlight that multi-agent systems can consist of manifold components (e.g., communication structure, agent profiles, memory).
%They discuss that agents acquire feedback from the environment, other agents, or human intervention. \citep{YinSCG23a, DuLTT23}

% Praising LLM Agents Research
\noindent{\textbf{Supportive Works.}}
% The potential benefits of MAD are increasingly recognized within recent literature.
Several works highlight the strengths of MAD, including divergent thinking \citep{degenofthought}, reasoning \citep{YinSCG23a}, creative writing \citep{SchickDJP22a}, dialogue generation \citep{ChenSB24a}, and theory of mind \citep{LiCSC23a}.
These works often use prompted personas \citep{WangMWG23a, XuYLW23a} and self-correction mechanisms like self-consistency \citep{WangWSL23} in a conversational setup.
However, the heterogeneous experimental setup across these studies (e.g., agent orchestration, decision-making, prompting) hinders their comparability \citep{Becker24, GuoCWC24a}.

% Critical LLM Agents Research
\noindent{\textbf{Critical Works.}}
Other researchers study the limitations of multi-agent systems.
Work led by Microsoft and Salesforce recently showed that many LLMs perform worse in multi-turn settings on generative problems, suffering from a systematic performance degradation \citep{laban2025llmslost}.
% \footnote{We note that our work is a resubmission from Feb 2025, while the mentioned paper was published in May 2025.}
\citet{WangWST24} and \citet{zhang2025multiagentdebateanswerquestion} focus on the computational cost of MAD, showing that single-agent LLMs can often achieve similar or even better performance through prompting.
% \citet{YinSCG23a} highlight the computational cost of single-model and multi-agent systems that try to improve reasoning.
Systems that include a self-critique mechanism (e.g., MAD) might diminish planning performance \citep{ValmeekamMK23a}. 
The correctness and content of LLM-generated criticism can be irrelevant to the performance of iterative prompting \citep{StechlyMK23a}.
\citet{ShahW24} argue that MAD has many issues, such as generalization, scalability, coordination, robustness, and ethical concerns.
% While multi-agent systems powered by LLMs show a prominent direction in solving complex tasks, this paradigm also comes with its costs and critiques. %make this back when camera ready

% Where do we step in?
\noindent{\textbf{Research Gap.}}
We highlight that MAD is relevant for problems that require divergent input or modular solutions.
While existing literature highlights diminished returns with prolonged MAD, the causes for such remain underexplored.
%While these systems offer potential advantages, their limitations require further systematic investigation.
We address this gap by characterizing the deterioration of MAD performance (RQ1), its prevalence (RQ2), and whether agents can recover from it (RQ3).
We systematically study possible causes and effects through a human evaluation (RQ4) and propose a first baseline to detect (RQ5) and recover (RQ6) the degrading discussions.
%By understanding these characteristics, we aim to uncover the underlying reasons and improve the design of such systems.